\documentclass[11pt]{article}
\usepackage[margin=1in]{geometry}
\usepackage{amsmath,amssymb,mathtools}
\usepackage{hyperref}
\usepackage{microtype}
\usepackage[numbers]{natbib}

\title{SCORES: Shelf Trajectories with Slow and Fast Components via Spectral Quadratic Surrogates \\ and Data Determined Weighting}
\author{}
\date{}

\begin{document}
\maketitle

\begin{abstract}
We study online reordering of a finite set of items to reduce user navigation time while remaining consistent with a drifting physical shelf layout.
Swipe adjacency is a policy artifact because the current order shapes the swipe path, so we learn proximity only from consecutive update events.
To obtain a solver with clean spectral structure, we replace absolute distance costs by a Laplacian quadratic surrogate that diagonalizes in the graph Fourier basis.
We then formulate a two time scale trajectory problem, with a slow shelf component and a fast app component coupled by a proximity term and controlled by temporal smoothness.
The key modification in this paper is that all weights and time scales are determined from the data: similarities are filtered by an adaptive state space estimator, coupling and smoothness weights are inverse variance weights, and the shelf low pass strength is selected by generalized cross validation.
A practical solution is a pipeline: spectral embeddings for daily structure, variance weighted temporal smoothing for shelf inference, then a noise floor stopped local swap search that outputs a valid permutation with stability emerging from uncertainty rather than manual caps.
\end{abstract}

\paragraph{Keywords.}
Online reordering; spectral seriation; Laplacian quadratic surrogate; time frequency smoothing; drifting shelf inference; adaptive filtering; generalized cross validation.

\tableofcontents
\bigskip

\section{Setup}
\subsection{Items and permutations}
Let $[n]:=\{1,\dots,n\}$ index items.
An ordering is a permutation $\pi$ of $[n]$.
Write $\mathrm{pos}_{\pi}(i)\in[n]$ for the position of item $i$ under $\pi$.

\subsection{Sessions, updates, and why swipe adjacency is not the right signal}
A session is a sequence of visited items $(v_1,\dots,v_m)$ with action indicators $a_t\in\{0,1\}$ where $a_t=1$ denotes an update event.
The swipe path $(v_t)$ depends strongly on the current UI order, so transition counts between consecutive swipes largely reflect current adjacency rather than user intent.

We therefore base learning on the subsequence of consecutive updates.
Let $t_1<\cdots<t_k$ be indices with $a_{t_r}=1$.
Define the update sequence
\[
u_r := v_{t_r},\qquad r=1,\dots,k.
\]
Optionally, collapse consecutive repeats by removing $u_{r+1}$ whenever $u_{r+1}=u_r$.
Let $\Delta_r\ge 0$ be the time between update events $r$ and $r+1$.

\subsection{Update adjacency weights with no clip constants}
We weight consecutive update pairs by a robust, scale free function of inter update time.
Let $\widehat{F}_d$ be the empirical cumulative distribution function of $\{\Delta_r\}$ pooled across sessions on day $d$.
Define
\[
w_d(\Delta):=1-\widehat{F}_d(\Delta)\in[0,1].
\]
This uses only ranks within the day, so it is invariant to absolute time scaling and introduces no clip threshold.
Other equivalent choices include $w_d(\Delta)=\mathrm{rank}(\Delta)/(\#\Delta)$.

For a day index $d$, define symmetric update adjacency weights
\[
\widehat{S}_{ij}(d)
:=
\sum_{\text{sessions on day }d}\ \sum_{r=1}^{k-1}
\mathbf{1}\{u_r=i,\ u_{r+1}=j\}\,w_d(\Delta_r)
\;+\;
\mathbf{1}\{u_r=j,\ u_{r+1}=i\}\,w_d(\Delta_r).
\]
This captures proximity only between items that are edited consecutively, which is exactly the quantity the UI order can realistically reduce.

\subsection{Adaptive two time scale evidence via state space filtering}
Fast and slow exponentially weighted averages require forgetting factors.
We use an adaptive state space estimator that produces a slow latent similarity and a fast residual without choosing half lives by hand.

For each unordered pair $(i,j)$ that appears with nonzero mass, let the observed daily similarity be $\widehat{S}_{ij}(d)$.
Model a latent slow similarity $S_{ij}(d)$ as a random walk observed in noise:
\begin{align}
\widehat{S}_{ij}(d) &= S_{ij}(d) + \varepsilon_{ij}(d), \\
S_{ij}(d) &= S_{ij}(d-1) + \xi_{ij}(d),
\end{align}
where $\varepsilon_{ij}(d)$ is observation noise and $\xi_{ij}(d)$ is process drift.

Define the fast (innovation/residual) similarity
\[
S^{\mathrm{fast}}_{ij}(d):=\widehat{S}_{ij}(d)-S_{ij}(d).
\]
Residual similarities can be negative. A standard Laplacian $L=D-S$ assumes nonnegative similarities if one wants a positive semidefinite quadratic energy and stable spectral structure; with negative weights, the Laplacian quadratic form can become indefinite and the Fiedler-vector interpretation breaks.
In practice, when building a ``fast'' Laplacian (and when scoring candidate swaps), use only the positive part
\[
S^{\mathrm{fast},+}_{ij}(d):=\max\{S^{\mathrm{fast}}_{ij}(d),0\}.
\]

Residual similarities can be negative. A standard Laplacian $L=D-S$ assumes nonnegative similarities if one wants a positive semidefinite quadratic energy and stable spectral structure; with negative weights, the Laplacian quadratic form can become indefinite and the Fiedler-vector interpretation breaks.
In practice, when building a ``fast'' Laplacian (and when scoring candidate swaps), use only the positive part
\[
S^{\mathrm{fast},+}_{ij}(d):=\max\{S^{\mathrm{fast}}_{ij}(d),0\}.
\]

A scalar Kalman filter per edge yields the update
\begin{align}
\text{Predict:}\quad &\widetilde{S}_{ij}(d) = S_{ij}(d-1), \\
\text{Innovation:}\quad &\nu_{ij}(d) = \widehat{S}_{ij}(d)-\widetilde{S}_{ij}(d), \\
\text{Gain:}\quad &K_{ij}(d) = \frac{P_{ij}(d-1)+Q_{ij}(d)}{P_{ij}(d-1)+Q_{ij}(d)+R_{ij}(d)}, \\
\text{Update:}\quad &S_{ij}(d) = \widetilde{S}_{ij}(d) + K_{ij}(d)\,\nu_{ij}(d), \\
\text{Covariance:}\quad &P_{ij}(d) = \bigl(1-K_{ij}(d)\bigr)\bigl(P_{ij}(d-1)+Q_{ij}(d)\bigr).
\end{align}

The remaining question is how to set the variances $R_{ij}(d)$ and $Q_{ij}(d)$ without hyperparameters.
We estimate them online from the data using the innovation sequence.
A simple, fully data determined option is to use rolling empirical second moments computed on the innovation stream itself:
\begin{align}
R_{ij}(d) &:= \mathrm{Var}\bigl(\widehat{S}_{ij}(d)-\widehat{S}_{ij}(d-1)\bigr), \\
Q_{ij}(d) &:= \mathrm{Var}\bigl(S_{ij}(d)-S_{ij}(d-1)\bigr),
\end{align}
where $\mathrm{Var}(\cdot)$ denotes an empirical variance computed over a trailing window that is itself selected by predictive likelihood, as described below.

To avoid choosing the window length, we select it by maximizing one step ahead predictive log likelihood on held out days:
\[
L^\star \in \arg\max_{L\in\mathcal{L}} \sum_{d\in\mathcal{D}_{\text{val}}} \log p\bigl(\widehat{S}_{ij}(d)\mid \widehat{S}_{ij}(d-1),\dots,\widehat{S}_{ij}(d-L)\bigr),
\]
where $\mathcal{L}$ is a small fixed candidate set such as $\{3,5,7,14,21,30\}$ days.
This is not a tunable parameter in deployment because it is selected by validation and can be reselected periodically.

We define:
\[
S^{\text{slow}}(d) := S(d),
\qquad
S^{\text{fast}}(d) := \widehat{S}(d) - S(d).
\]
The slow component captures persistent structure, the fast component captures day specific deviations.

\subsection{Cold item uncertainty from data}
Cold item protection emerges from uncertainty.

Let $\widehat{c}_i(d)$ be the number of update events involving item $i$ on day $d$:
\[
\widehat{c}_i(d) := \sum_{\text{sessions on day }d}\sum_{r=1}^{k}\mathbf{1}\{u_r=i\}.
\]
Define an activity rate model with Poisson observations and a drifting log rate:
\begin{align}
\widehat{c}_i(d) &\sim \mathrm{Poisson}\bigl(\lambda_i(d)\bigr), \\
\log \lambda_i(d) &= \log \lambda_i(d-1) + \zeta_i(d).
\end{align}
A one dimensional state space filter yields an estimate of $\lambda_i(d)$ and its posterior variance.
Items with low activity have high uncertainty, which will translate into higher stability weights downstream.

\section{Quadratic surrogate that diagonalizes on graphs}
\subsection{Laplacians and scale normalization}
Given a symmetric similarity matrix $S(d)$, define degree $D(d)=\mathrm{diag}(S(d)\mathbf{1})$ and Laplacian
\[
L(d) := D(d) - S(d).
\]
For an embedding vector $x\in\mathbb{R}^n$,
\[
x^\top L(d) x = \frac12\sum_{i=1}^n\sum_{j=1}^n S_{ij}(d)\,(x_i-x_j)^2.
\]
This is the squared distance surrogate. It is quadratic and diagonalizes in the eigenbasis of $L(d)$.

Daily traffic can scale $S(d)$ up or down.
To eliminate scale sensitivity, we normalize Laplacians by trace:
\[
\overline{L}(d) := \frac{L(d)}{\mathrm{tr}(L(d))+\mathbf{1}\{\mathrm{tr}(L(d))=0\}},
\]
so that energies are comparable across days.
For the slow and fast components we use
\[
\overline{L}^{\text{slow}}(d)\ \text{from}\ S^{\text{slow}}(d),
\qquad
\overline{L}^{\text{fast}}(d)\ \text{from}\ S^{\text{fast}}(d).
\]

\subsection{From embeddings to permutations}
Given an embedding $x$, define the induced ordering
\[
\pi(x) := \mathrm{argsort}(x),
\]
sorting items by $x_i$.

Spectral seriation uses the Fiedler vector, an eigenvector associated with the second smallest eigenvalue of $\overline{L}(d)$:
\[
x^{\star}(d)\in\arg\min_{x\perp\mathbf{1},\ \|x\|_2=1} x^\top \overline{L}(d) x.
\]

\paragraph{Production stability: sign ambiguity and eigengap degeneracy.}
In production, eigenvector-based orderings can ``thrash'' even when the underlying graph changes little.
Three common causes are: (i) the arbitrary sign of eigenvectors, (ii) near-degeneracy when $\lambda_2$ is close to $\lambda_3$ (the Fiedler direction is poorly identified), and (iii) changes in graph connectivity on sparse days.
These manifest as day-to-day full reversals (sign flips), block-level jitter (rotation within a nearly-degenerate eigenspace), and meaningless embeddings when the graph is too disconnected.

\paragraph{Fix:}
After computing $x^\star(d)$, enforce alignment with the previous day:
\[
\langle x^\star(d),x^\star(d-1)\rangle \ge 0.
\]
If $\langle x^\star(d),x^\star(d-1)\rangle < 0$, replace $x^\star(d)\leftarrow -x^\star(d)$.
This removes global day-to-day reversals due purely to sign ambiguity.

\section{Trajectory formulation with slow shelf and fast app components}
\subsection{State variables}
For each day $d\in\{0,\dots,D-1\}$:
\begin{enumerate}
\item $x_a(d)\in\mathbb{R}^n$ is the fast app embedding.
\item $x_s(d)\in\mathbb{R}^n$ is the slow shelf embedding.
\item $\pi_a(d)=\pi(x_a(d))$ is the shown order after rounding and constraints.
\item $\pi_s(d)=\pi(x_s(d))$ is the inferred shelf order after rounding.
\end{enumerate}

\subsection{Data terms}
We use the Laplacian quadratic surrogate for both components:
\begin{equation}
J_{\text{graph}} :=
\sum_{d=0}^{D-1} x_a(d)^\top \overline{L}^{\text{fast}}(d)\,x_a(d)
\;+\;
\sum_{d=0}^{D-1} x_s(d)^\top \overline{L}^{\text{slow}}(d)\,x_s(d).
\label{eq:data_terms_hplfree}
\end{equation}

\subsection{Variance weighted coupling and shelf smoothness}
We propose a quadratic trajectory objective:
\begin{align}
\min_{\{x_a(d),x_s(d)\}}
\ &J_{\text{graph}}
\;+\;
\sum_{d=0}^{D-1} \|x_a(d)-x_s(d)\|_{\Gamma(d)}^2
\label{eq:coupling_hplfree}
\\
&\;+\;
\sum_{d=1}^{D-1} \|x_s(d)-x_s(d-1)\|_{K(d)}^2
\;+\;
\sum_{d=2}^{D-1} \|x_s(d)-2x_s(d-1)+x_s(d-2)\|_{H(d)}^2,
\label{eq:temporal_smooth_hplfree}
\end{align}
with constraints per day:
\[
x_a(d)\perp\mathbf{1},\ \|x_a(d)\|_2=1,
\qquad
x_s(d)\perp\mathbf{1},\ \|x_s(d)\|_2=1.
\]

Here $\|z\|_{M}^2 := z^\top M z$ with diagonal positive semidefinite matrices:
\begin{enumerate}
\item $\Gamma(d)$ couples app and shelf.
\item $K(d)$ penalizes shelf velocity.
\item $H(d)$ penalizes shelf acceleration.
\end{enumerate}

All three are determined by inverse variance weighting.
Let $R_i(d)$ be an estimate of $\mathrm{Var}\bigl(x_{a,i}(d)-x_{s,i}(d)\bigr)$.
Let $Q_{1,i}(d)$ estimate $\mathrm{Var}\bigl(x_{s,i}(d)-x_{s,i}(d-1)\bigr)$.
Let $Q_{2,i}(d)$ estimate $\mathrm{Var}\bigl(x_{s,i}(d)-2x_{s,i}(d-1)+x_{s,i}(d-2)\bigr)$.
Then set
\[
\Gamma_{ii}(d) := \frac{1}{R_i(d)},
\qquad
K_{ii}(d) := \frac{1}{Q_{1,i}(d)},
\qquad
H_{ii}(d) := \frac{1}{Q_{2,i}(d)}.
\]
Cold item protection is automatic because low activity implies high uncertainty in the inferred shelf coordinates, which increases the effective weights and resists motion.

\section{Frequency domain view of shelf drift and data selected low pass strength}
This section explains the slow and fast split in time frequency terms.
For clarity, consider the shelf smoothing subproblem with $x_a(d)$ fixed:
\begin{equation}
\min_{\{x_s(d)\}}
\ \sum_{d=0}^{D-1}\|x_a(d)-x_s(d)\|_{\Gamma}^2
+
\sum_{d=1}^{D-1}\|x_s(d)-x_s(d-1)\|_{K}^2
+
\sum_{d=2}^{D-1}\|x_s(d)-2x_s(d-1)+x_s(d-2)\|_{H}^2,
\label{eq:shelf_smooth_only_hplfree}
\end{equation}
with constant diagonal weights for exposition.

Let $X_a(\omega)$ and $X_s(\omega)$ be the discrete Fourier transforms over $d$ applied coordinate wise.
Define
\[
g(\omega)=4\sin^2(\pi\omega/D).
\]
Then the first difference term contributes a factor $g(\omega)$ and the second difference contributes $g(\omega)^2.
$
The objective diagonalizes by $\omega$ and by coordinate:
\[
\sum_{\omega=0}^{D-1}
\left(
\|X_a(\omega)-X_s(\omega)\|_{\Gamma}^2
+
g(\omega)\|X_s(\omega)\|_{K}^2
+
g(\omega)^2\|X_s(\omega)\|_{H}^2
\right).
\]
The minimizer has the explicit low pass form, coordinate wise:
\[
X_{s,i}(\omega) =
\frac{\Gamma_{ii}}{\Gamma_{ii}+K_{ii}\,g(\omega)+H_{ii}\,g(\omega)^2}\ X_{a,i}(\omega).
\]
Low frequencies with small $g(\omega)$ pass into the shelf, high frequencies are suppressed.
The residual $X_a(\omega)-X_s(\omega)$ is a complementary high pass, which is the fast component injected into the UI.

The remaining issue is selecting the effective low pass strength without manual ratios.
We reparameterize by tying scales to $\Gamma$:
\[
K=\alpha\,\Gamma,
\qquad
H=\beta\,\Gamma,
\]
so the gain depends only on $(\alpha,\beta)$.
We choose $(\alpha,\beta)$ by generalized cross validation.
Let $\widehat{x}_s(\alpha,\beta)=A_{\alpha,\beta}x_a$ denote the linear smoothing operator in time.
Then define
\[
\mathrm{GCV}(\alpha,\beta):=
\frac{\|x_a-\widehat{x}_s(\alpha,\beta)\|_2^2}{\bigl(1-\mathrm{tr}(A_{\alpha,\beta})/D\bigr)^2}.
\]
We select
\[
(\alpha^\star,\beta^\star)\in\arg\min_{(\alpha,\beta)\in\mathcal{G}} \mathrm{GCV}(\alpha,\beta),
\]
over a fixed grid $\mathcal{G}$ on log scale.
This selection is data determined and can be recomputed periodically.

\section{Practical solver}
\subsection{Stage 1: build observed similarities from update events}
On each day $d$:
\begin{enumerate}
\item Parse sessions to obtain update sequences $(u_r)$ and inter update times $\Delta_r$.
\item Compute $w_d(\Delta)=\widehat{F}_d(\Delta)$ from that day distribution.
\item Accumulate $\widehat{S}(d)$ via consecutive updates.
\item Update per edge latent slow similarity $S(d)$ with the adaptive state space filter.
\item Set $S^{\text{slow}}(d)=S(d)$ and $S^{\text{fast}}(d)=\widehat{S}(d)-S(d)$.
\item Form normalized Laplacians $\overline{L}^{\text{slow}}(d)$ and $\overline{L}^{\text{fast}}(d)$.
\end{enumerate}

\subsection{Stage 2: spectral embeddings}
Compute daily fast embedding $x_a^{(0)}(d)$ as the Fiedler vector of $\overline{L}^{\text{fast}}(d)$.
Compute daily slow embedding $x_s^{(0)}(d)$ as the Fiedler vector of $\overline{L}^{\text{slow}}(d)$.

These are graph frequency solutions, diagonalized by the eigenbasis of each Laplacian.

\subsection{Stage 3: infer a drifting shelf by variance weighted time smoothing}
Use \eqref{eq:shelf_smooth_only_hplfree} to smooth the slow embedding trajectory over time.
Two equivalent implementations:
\begin{enumerate}
\item Frequency domain: take the DFT over $d$, apply the low pass gain coordinate wise, then invert the DFT.
\item Time domain: solve the banded normal equations induced by \eqref{eq:shelf_smooth_only_hplfree}, separately for each coordinate.
\end{enumerate}
Set $K=\alpha^\star\Gamma$ and $H=\beta^\star\Gamma$ with $(\alpha^\star,\beta^\star)$ chosen by GCV.
Denote the smoothed result by $x_s(d)$ and set $\pi_s(d)=\pi(x_s(d))$.

\subsection{Stage 4: inject fast improvements onto the shelf using a noise floor stopping rule}
We produce a valid shown permutation by starting from the shelf order and applying local moves guided by the fast quadratic surrogate.
Let $\pi\gets\pi_s(d)$.

Define the discrete fast surrogate consistent with the Laplacian energy:
\[
\widetilde{J}_{\text{fast}}(\pi;d) := \sum_{i<j} S^{\text{fast}}_{ij}(d)\ \big(\mathrm{pos}_{\pi}(i)-\mathrm{pos}_{\pi}(j)\big)^2.
\]
This is the squared distance surrogate, so it matches the diagonalizable form used for embeddings.

\paragraph{Candidate move set without window parameters.}
Rather than restricting swaps by position radius, we restrict by graph support.
For each item $i$, let $\mathcal{N}_d(i)$ be the set of neighbors with the largest weights $S^{\text{fast}}_{ij}(d)$.
We choose the size of $\mathcal{N}_d(i)$ by a data defined sparsification rule, for example keeping all neighbors with $S^{\text{fast}}_{ij}(d)$ above the day median of nonzero weights.
This introduces no manual $k$.

Candidate swaps are pairs $(i,j)$ with $j\in\mathcal{N}_d(i)$.

\paragraph{Noise floor estimate.}
Compute a noise scale for swap improvements by sampling a modest number of random candidate swaps from the same candidate set and evaluating their $\Delta \widetilde{J}_{\text{fast}}$ under the current $\pi$.
Let $\widehat{\sigma}_{\Delta J}(d)$ be the empirical standard deviation of these sampled improvements.

\paragraph{Greedy search with automatic stopping.}
Iterate:
\begin{enumerate}
\item Find the candidate swap with the most negative improvement $\Delta \widetilde{J}_{\text{fast}}$.
\item Accept it if $\Delta \widetilde{J}_{\text{fast}} \le -\widehat{\sigma}_{\Delta J}(d)$ and feasibility constraints are satisfied.
\item Stop when no candidate satisfies the noise floor threshold.
\end{enumerate}

This yields a natural disruption level determined by signal to noise.
On low signal days the algorithm stops quickly and preserves the shelf order.
On strong signal days it performs more swaps, but only when improvements are clearly above noise.

\paragraph{Stability emerges from uncertainty rather than caps.}
Cold items rarely have strong edges in $S^{\text{fast}}(d)$, so they are rarely proposed as beneficial swaps.
They also have higher uncertainty in shelf inference, which increases coupling and smoothness weights and makes their inferred shelf positions more stable.
No manual displacement caps are required.

Return $\pi_a(d)=\pi$.

\subsection{Complexity}
For sparse similarities the Laplacian is sparse.
Computing a Fiedler vector can be done with iterative methods whose dominant cost is sparse matrix vector products.
The time smoothing step solves banded systems or applies FFTs over $d$.
The swap phase evaluates a sparse set of candidate swaps defined by fast graph neighbors, and can be implemented with incremental score updates that touch only edges incident to swapped items.

\paragraph{Big-$O$ summary.}
Let $n$ be the number of items, $D$ the number of days, and let $m_d:=\#\{(i,j): S_{ij}(d)\neq 0\}$ be the number of nonzero similarities (edges) on day $d$ with $m:=\max_d m_d$.
\begin{enumerate}
\item \emph{Daily spectral embedding.} Using Lanczos/power iterations, $T$ iterations of sparse matrix--vector multiplies cost
\[
O\bigl(T\,m_d\bigr)\ \text{per day, hence}\ O\bigl(T\sum_{d=0}^{D-1} m_d\bigr)=O\bigl(T D m\bigr).
\]
\item \emph{Temporal smoothing (per coordinate).} If the shelf smoothing is implemented as a banded linear solve with constant bandwidth, it costs $O(D)$ per coordinate, hence $O(nD)$ total. If implemented via FFT-based convolution in time, it costs $O(D\log D)$ per coordinate, hence $O(nD\log D)$ total.
\item \emph{Local swap search.} If each item proposes $k$ candidate neighbors (e.g., from the fast graph), the candidate set size is $O(nk)$; with incremental updates that touch only incident edges, each evaluation/update is $O(k)$, so $I$ accepted iterations cost $O(I\,k)$ plus an initial pass $O(nk)$.
\end{enumerate}

\paragraph{End-to-end complexity (SCORES).}
Under the assumptions above and ignoring lower-order terms, the overall cost across $D$ days is
\[
O\bigl(T D m\bigr)\;+\;
\begin{cases}
O(nD), & \text{banded time-smoothing},\\
O(nD\log D), & \text{FFT-based time-smoothing},
\end{cases}
\;+\; O(nk + I k).
\]

\section{Discussion: what is now data determined}
SCORES has three places where hyperparameters typically enter: time scale, stability weights, and swap disruption controls.
In this formulation:
\begin{enumerate}
\item Time scale is determined by an adaptive state space estimator, with effective memory governed by estimated process and observation variance.
\item Stability weights $\Gamma,K,H$ are inverse variance weights derived from uncertainty in coupling residuals and shelf differences.
\item Swap disruption is controlled by a noise floor stopping rule, so the number of swaps is an output, not an input.
\end{enumerate}

\section{Conclusion}
SCORES is formulated as a two time scale trajectory problem.
Update to update adjacency provides an intent aligned signal and avoids self reinforcing swipe adjacency.
A Laplacian squared distance surrogate yields spectral diagonalization over items, and temporal difference penalties yield diagonalization over time frequencies.
In this paper, the time scales, weights, and disruption controls are selected from the data through adaptive filtering, inverse variance weighting, and generalized cross validation.
The resulting system supports a drifting shelf estimate and bounded fast injections into the shown ordering with stability emerging from uncertainty rather than manual tuning.

\bibliographystyle{plainnat}
\end{document}
